\documentclass[10pt,twocolumn]{article}

\usepackage[a4paper,top=0.9in,bottom=0.9in,left=0.6in,right=0.6in,columnsep=0.2in]{geometry}
\usepackage{times}
\usepackage{graphicx}
\usepackage{booktabs}
\usepackage{multirow}
\usepackage{bm}
\usepackage{microtype}
\usepackage{hyperref}
\usepackage{listings}
\usepackage{xcolor}
\usepackage{amsmath}
\usepackage{caption}
\usepackage{titlesec}

\hypersetup{colorlinks=true, linkcolor=blue, citecolor=blue, urlcolor=blue}

\titleformat{\section}{\normalfont\large\bfseries}{\thesection}{0.5em}{}
\titleformat{\subsection}{\normalfont\normalsize\bfseries}{\thesubsection}{0.5em}{}
\titlespacing*{\section}{0pt}{1.2ex}{0.6ex}
\titlespacing*{\subsection}{0pt}{0.8ex}{0.4ex}

\setlength{\parskip}{2pt}
\setlength{\parindent}{1em}

\lstset{
  basicstyle=\ttfamily\footnotesize,
  breaklines=true,
  columns=fullflexible,
  keywordstyle=\color{blue!70!black}\bfseries,
  commentstyle=\color{green!50!black}\itshape,
  stringstyle=\color{red!50!black},
  numbers=none,
  frame=single,
  framesep=4pt,
  language=C++,
  escapeinside={(*@}{@*)}
}

\title{\vspace{-2em}\textbf{Throughput-Probed Bandit (TPB): Zero-Hot-Path-Overhead\\
  Adaptive Promotion for Predictive-Translation Buffer Pools}}
\author{
  \textit{Anonymous Author(s)}\\
  \textit{Reproduction-with-extension report on PrediCache (SIGMOD 2026)}
}
\date{}

\begin{document}
\twocolumn[
  \maketitle
  \begin{abstract}
    \noindent
    PrediCache (SIGMOD 2026) introduces \emph{predictive translation}: each
    page is bound to a deterministic preferred buffer frame, letting the CPU
    speculatively read from that frame in parallel with the hash-table lookup.
    Accuracy is maintained by probabilistically promoting hot pages into their
    preferred frames with hand-picked constants ($\tfrac{1}{32}$ for the
    no-displace case, $\tfrac{1}{512}$ otherwise).
    We show empirically that these constants are workload-dependent and
    introduce \textbf{TPB}, an online controller that adapts the probability
    using the system's own per-second transaction counter as its only signal.
    TPB's key architectural property is \emph{zero hot-path overhead}: a
    single global atomic-relaxed load per buffer access, identical to the
    static default.  All controller logic lives on the existing stat thread.
    Across TPC-C, uniform random-read, and skewed YCSB
    ($\zeta\!=\!0.9$), TPB matches the hand-picked default within
    statistical noise without manual per-workload tuning, eliminating
    a 22\,\% regression that an earlier non-architectural variant suffered on
    skewed reads.  We also document the v3$\to$v4$\to$v5 development arc:
    the lessons learned about \emph{where} to measure are
    more general than TPB itself.
  \end{abstract}
  \vspace{1.5em}
]

% ======================================================================
\section{Introduction}
% ======================================================================

Buffer-pool design has been driven for a decade by a tight efficiency target:
each page-id (PID)-to-frame translation should cost no more than the
in-memory case~\cite{leis2018leanstore,leis2023vmcache}.  The
hash-table-based pools of MySQL, PostgreSQL, and SQLite carry two
\emph{dependent} memory accesses per page request (chained look-up plus
header dereference), which has historically been seen as a fundamental
bottleneck~\cite{leis2018leanstore,leis2023vmcache}.  This motivated
pointer-swizzling (LeanStore~\cite{leis2018leanstore}) and OS-page-table
designs (vmcache~\cite{leis2023vmcache}), each of which sacrifices a
qualitative property in exchange for speed.

\textbf{PrediCache}~\cite{zinsmeister2026predictive} closes this gap without
either trade-off.  Every PID is bound at \emph{startup} to a deterministic
preferred frame ($\textsf{hash}(\textsf{pid}) \bmod \textsf{frameCount}$).
On a page request, the CPU speculatively reads from the predicted frame
\emph{in parallel} with the hash-table look-up; superscalar execution hides
one of the two memory accesses.  When the prediction is correct, two
dependent accesses are completed in the time of one.  Accuracy of the
speculation is maintained by \emph{probabilistic promotion}: on the second
and later accesses to a page, it is moved to its preferred frame with
fixed probability $p_0 = \tfrac{1}{32}$, or with probability
$p_1 = \tfrac{1}{512}$ when the move requires displacing an incumbent
page~\cite{zinsmeister2026predictive}.

These two constants are the only knobs in PrediCache that are
\emph{not justified} by a first-principles argument in the paper.  Section~3
shows empirically that $p_0\!=\!\tfrac{1}{32}$ is suboptimal by 1--2\,\%
on TPC-C and uniform random-read, and exactly optimal on
skewed YCSB at $\zeta\!=\!0.9$.  In other words,
\textbf{the right $p_0$ is workload-dependent}.

A naïve adaptation strategy --- per-access counters, hill-climbing on
hit-rate --- fails catastrophically in this regime.  Our earlier attempts
(CBP-v3, CBP-v4, TPB-v1\,\dots v3) ranged from $-22.4\,\%$ to $-2.2\,\%$
versus the static default on a 30-second TPC-C run; the worst case was
significant at $|t|\!=\!4.48$.  Two failure modes dominated: \emph{(i)}
per-access bookkeeping (atomic counter increments, sampling) imposes
8--23\,\% throughput tax even when no adaptation runs; \emph{(ii)} hit-rate
is monotonic in $\log_2(1/p_0)$, so any bandit that uses hit-rate as its
reward overshoots into one of the extremes (always-promote or
never-promote) rather than stopping at the throughput optimum.

\textbf{Contribution of this paper.}
We design and evaluate \textbf{TPB} (Throughput-Probed Bandit), an
adaptive promoter that solves both failure modes by a single architectural
choice: \emph{the instrumentation lives outside the buffer-pool access path
entirely}.  TPB uses a three-window probe-sandwich
$[\text{lc}_{\text{base}}, \text{lc}_{\text{probe}}, \text{lc}_{\text{base}}]$
fed by PrediCache's own per-second stat thread; the entire controller
state machine is $\approx 50$ lines of C++ and runs once per second on a
single thread.  Hot-path access cost is one relaxed atomic load per page
request --- identical to the unmodified static default.

We make three concrete claims:

\begin{enumerate}\itemsep0pt
\item TPB matches PrediCache's static $\tfrac{1}{32}$ default within
      one standard deviation on all three workloads tested (TPC-C,
      Random Read, Skewed YCSB) at $n\!=\!5$ replicates.
\item TPB \emph{eliminates} the catastrophic 22.4\,\% regression seen
      with hot-path-instrumented controllers on Skewed YCSB
      (26\,M TX/s class), because the instrumentation is now outside
      the per-access path.
\item The development arc (v3$\to$v4$\to$v5) yields a general lesson
      for any online-control problem in a high-throughput data-plane
      system: \emph{move instrumentation to the system's existing
      coarse-grained boundaries} before designing the controller.
\end{enumerate}

Sections 2--3 review PrediCache and present the sensitivity finding that
motivates adaptation.  Section 4 details TPB.  Section 5 evaluates against
the static default.  Section 6 covers related work.  Section 7 lists
limitations.

% ======================================================================
\section{Background: PrediCache and its Promotion Lifecycle}
% ======================================================================

Pages in PrediCache live in one of two states: the \emph{preferred} frame
(deterministically computed from the PID) or an \emph{arbitrary} frame
elsewhere in the buffer pool.  Reading a page in its preferred frame
takes one effective memory access (the hash-table look-up runs in parallel
under superscalar execution).  Reading it in an arbitrary frame takes
two dependent memory accesses.

The page \emph{lifecycle} is:

\begin{enumerate}\itemsep0pt
\item \textbf{First access:} the page is loaded from storage into an
      arbitrary free frame.  PrediCache assumes the page may be a
      ``one-hit wonder'' and avoids the cost of immediately placing it
      at its preferred slot~\cite{zinsmeister2026predictive}.
\item \textbf{Subsequent access:} with probability $p_0\!=\!\tfrac{1}{32}$,
      the page is \emph{promoted} (\texttt{memcpy} of 4\,KB plus two
      exclusive latches) to its preferred frame.  If the preferred
      frame is occupied by a hotter page, a $\tfrac{1}{16}$ probability
      gate is applied on top, yielding effective
      $p_1 = p_0 \cdot \tfrac{1}{16} = \tfrac{1}{512}$ for the displace
      case.
\item \textbf{Eviction:} on memory pressure, the page is dropped from
      whatever frame it occupies and the frame is returned to the free
      list.
\end{enumerate}

\noindent
Two observations about this design:

\textbf{(a) $p_0$ trades off two costs.}  Higher $p_0$ moves hot pages
into preferred frames faster, accelerating the rate at which the
fast-path applies to subsequent accesses.  Lower $p_0$ avoids the
4\,KB \texttt{memcpy} bandwidth.  Neither extreme is universally optimal.

\textbf{(b) $p_0$ has no closed-form best value.}  The fast-path benefit
per promotion depends on a page's expected residency and reuse rate,
which is workload-dependent.  PrediCache picks $p_0 = \tfrac{1}{32}$
without a derivation; the paper simply states ``we
chose''~\cite[\S5.1]{zinsmeister2026predictive}.

% ======================================================================
\section{Sensitivity Analysis: $\tfrac{1}{32}$ is workload-dependent}
% ======================================================================

We sweep $p_0$ over $\{0,\,\tfrac{1}{16},\,\tfrac{1}{32},\,\tfrac{1}{128}\}$
on three workloads at 16 threads, in-memory.  $p_1$ tracks $p_0$
proportionally (see \S2).  Table~\ref{tab:sweep} summarizes the steady-state
mean throughput across two replicates.  Hardware and workload parameters
are documented in \S5.

\begin{table}[h]
\centering\small
\caption{Static-promotion sensitivity (n=2, 16 threads, in-memory).
  Bold = workload's argmax.}
\label{tab:sweep}
\begin{tabular}{lrrr}
\toprule
 $p_0$ & TPC-C (k TX/s) & Rnd-Read (k TX/s) & Skew-YCSB (M TX/s) \\
\midrule
$0$ (off)   & 633 ($-1.7\%$) & 15\,521 ($-1.7\%$) & 25.4 ($-4.4\%$) \\
$\tfrac{1}{16}$ & \textbf{651} ($+1.2\%$) & \textbf{16\,033} ($+1.6\%$) & 25.6 ($-3.7\%$) \\
$\tfrac{1}{32}$ & 644 (paper) & 15\,786 (paper) & \textbf{26.6} (paper) \\
$\tfrac{1}{128}$ & 622 ($-3.3\%$) & 16\,016 ($+1.5\%$) & 26.0 ($-2.0\%$) \\
\bottomrule
\end{tabular}
\end{table}

\noindent
Three structural findings:

\textbf{(F1) Concave curves on OLTP and uniform-random.}  Both
TPC-C and Random Read have a clear interior maximum at $p_0\!=\!\tfrac{1}{16}$.

\textbf{(F2) Optimum at $\tfrac{1}{32}$ for skewed reads.}
Skewed YCSB is monotonically worse on either side of $\tfrac{1}{32}$.
PrediCache's default \emph{is} the optimum for this workload class.

\textbf{(F3) Differences are small (1--3\,\%).}  The fixed-probability
curves are flat near the optimum.  This is a tight competitive surface
for any adaptive controller --- there is little headroom to win.

The TPC-C delta ($+1.2\,\%$ at $p_0\!=\!\tfrac{1}{16}$) is within one
$\sigma$ at $n\!=\!2$, so we mark F1 as \emph{suggestive}.  Direction is
consistent across replicates.

% ======================================================================
\section{TPB: Throughput-Probed Bandit}
% ======================================================================

% --------------------------------------------------------
\subsection{Design Goals}
% --------------------------------------------------------

We target three properties simultaneously:

\textbf{G1 (zero hot-path overhead).}  The per-access cost under TPB must
equal the cost under the unmodified static default.  Specifically: one
relaxed atomic load to read the current $\log_2(1/p_0)$ mask, and no
other operations.

\textbf{G2 (throughput-aligned signal).}  The controller's reward must
be the system's actual throughput, not a proxy like hit-rate that is
monotonic in $\log_2(1/p_0)$.

\textbf{G3 (no manual per-workload tuning).}  A single configuration
must match or beat the hand-picked default across all three of our
target workloads.

% --------------------------------------------------------
\subsection{Mechanism: Probe-Sandwich Bandit}
% --------------------------------------------------------

TPB maintains a single integer \texttt{lc\_base} (the current best
$\log_2(1/p_0)$) that all worker threads read on every access.  The
controller fires once per stat-thread window (1 s).  Every
\texttt{TPB\_PROBE\_EVERY} windows it opens a \emph{three-window sandwich}:

\[
\underbrace{\text{lc}_{\text{base}}}_{\text{window } k-1} \;\to\;
\underbrace{\text{lc}_{\text{probe}}}_{\text{window } k} \;\to\;
\underbrace{\text{lc}_{\text{base}}}_{\text{window } k+1}
\]

\noindent
where $\text{lc}_{\text{probe}} = \text{lc}_{\text{base}} \pm 1$, direction
alternating after each acceptance.  Let $T_{\text{pre}}, T_{\text{probe}},
T_{\text{post}}$ be the per-window transaction counts.  The controller:

\begin{enumerate}\itemsep0pt
\item Computes
  $T_{\text{avg}} = (T_{\text{pre}} + T_{\text{post}}) / 2$.
\item Requires
  $|T_{\text{pre}} - T_{\text{post}}| / T_{\text{pre}} < \tau_{\text{phase}}$.
  This phase guard rejects sandwiches where the workload was non-stationary
  during the probe (TPC-C's dataset grows over time, naturally causing a
  10--15\,\% per-30\,s decline).
\item Accepts the probe iff
  $T_{\text{probe}} > T_{\text{avg}} \cdot (1 + \tau_{\text{accept}})$.
\item On acceptance: $\text{lc}_{\text{base}} \leftarrow
  \text{lc}_{\text{probe}}$, flip probe-direction for the next sandwich.
\end{enumerate}

\noindent
Defaults: $\tau_{\text{phase}}\!=\!0.15$,
$\tau_{\text{accept}}\!=\!0.005$, \texttt{PROBE\_EVERY}=16.  The phase
guard tolerates the smooth monotonic drift of OLTP datasets while still
rejecting genuine non-linear phase changes.  The accept margin
($+0.5\,\%$) is the minimum throughput delta the controller will act on.

\textbf{Linear-interpolation rationale.}  $T_{\text{avg}}$ is the
linear interpolation of $\text{lc}_{\text{base}}$ throughput at the
probe's time slot.  Under the (weak) assumption that throughput drift
is linear over three seconds, $T_{\text{probe}} > T_{\text{avg}}$ is a
valid unbiased test of whether $\text{lc}_{\text{probe}}$ outperforms
$\text{lc}_{\text{base}}$ at that instant.

% --------------------------------------------------------
\subsection{Implementation}
% --------------------------------------------------------

The hot path reads a single global mask:

\begin{lstlisting}
inline uint64_t tpb::maskFor(unsigned cls) {
  uint8_t lc = tpb_lc_active.load(
      std::memory_order_relaxed);
  return (1u << lc) - 1u;
}
\end{lstlisting}

\noindent
This is identical to the static default except that the constant is
replaced by a single atomic-relaxed load --- on x86 a plain
\texttt{mov} from a hot cache line.  No fence, no
\texttt{compare\_exchange}, no per-class fan-out.  Per-access overhead
is \emph{algorithmically zero}.

The controller state machine is invoked from PrediCache's existing
per-second stat thread, which already drains
\texttt{txProgress.exchange(0)} for stdout logging:

\begin{lstlisting}
// Benchmark.hpp, inside the stat loop
u64 prog = txProgress.exchange(0);
cbp::tpbOnWindow(prog);   // (*@\bf{added: 1 line}@*)
\end{lstlisting}

\noindent
The full controller is $\approx 50$ lines of C++; the diff against
unmodified PrediCache is roughly 60 lines including environment-variable
parsing for the three knobs (\texttt{TPB\_INIT\_LC},
\texttt{TPB\_PHASE\_TOL}, \texttt{TPB\_PROBE\_EVERY}).

% --------------------------------------------------------
\subsection{Initial Direction and Boundary Handling}
% --------------------------------------------------------

A subtle correctness bug bit our earlier versions: the direction-update
\texttt{tpb\_last\_dir = (lc\_probe > lc\_base) ? +1 : -1} \emph{after}
the assignment \texttt{lc\_base = lc\_probe} compares equal values and
always evaluates to $-1$.  This produces a controller that always probes
``downward'' after each acceptance, causing degenerate convergence to
$p_0 = 1$ (always-promote).

The fix is to capture the acceptance direction \emph{before} mutating
\texttt{lc\_base}, and to flip the probe direction after each acceptance
so the controller samples both sides of the current best.  When the
candidate $\text{lc}_{\text{probe}}$ would step outside
$[\texttt{MIN\_LOG}, \texttt{MAX\_LOG}]\!=\![1,14]$, the direction is
inverted before computing the candidate.  This boundary-aware
direction-update is what enables the controller to escape the
catastrophic v3 regime described in \S5.5.

% ======================================================================
\section{Evaluation}
% ======================================================================

% --------------------------------------------------------
\subsection{Hardware and Reproduction Scope}
% --------------------------------------------------------

We run all experiments on a single dual-socket server: 2$\times$ Intel
Xeon Gold 5320 (52 cores, 1 thread/core, 2 NUMA nodes),
440\,GB DRAM, storage on a Lustre parallel filesystem with O\_DIRECT
support.  The PrediCache prototype~\cite{zinsmeister2026predictive_repo}
builds with \texttt{g++-11.4 -O3 -std=c++23 -laio}.

\textbf{Not reproduced.}  We cannot run the paper's
vmcache/LeanStore/WiredTiger/LMDB comparisons: vmcache requires the
\texttt{exmap} kernel module (security policy prohibits this on shared
infrastructure), and the other three engines require 1--3 days of
integration work outside the scope of this study.  All claims in this
paper are restricted to comparisons \emph{within} PrediCache.

\textbf{Workloads.}  All in-memory.

\begin{itemize}\itemsep0pt
\item \textbf{TPC-C:} 8 warehouses (initially $\approx 0.8$\,GB,
  growing to $\approx 9$\,GB over 30\,s), 16\,GB buffer pool.
\item \textbf{Random Read:} 10\,M uniform-distributed 128-byte
  key/value entries, 8\,GB buffer pool.
\item \textbf{Skewed YCSB:} 10\,M entries, Zipf $\zeta\!=\!0.9$,
  100\,\% read, 1 operation per transaction, 8\,GB buffer pool.
\end{itemize}

All runs are 30\,seconds, 16 worker threads, the first 5\,seconds
discarded as warmup, $n\!=\!5$ replicates per cell.  Numbers report
mean throughput in TX/s with sample standard deviation.

% --------------------------------------------------------
\subsection{TPB vs Static Default Across Three Workloads}
% --------------------------------------------------------

\begin{table*}[h]
\centering\small
\caption{TPB ($\texttt{INIT\_LC=5}, \texttt{PHASE\_TOL=15\%}, \texttt{PROBE\_EVERY=16}$)
  vs PrediCache's static $\tfrac{1}{32}$ default.  $n\!=\!5$, 30\,s, 16 threads.
  $|t|$ is Welch's two-sided statistic; significance at 95\,\% requires
  $|t|>2.31$.  Throughput is in thousands of transactions per second
  ($\sigma$ is sample standard deviation).}
\label{tab:tpb}
\setlength{\tabcolsep}{4pt}
\begin{tabular}{l l rrrr rr}
\toprule
Workload & Variant & Mean & $\sigma$ & min & max & $\Delta$ vs def & $|t|$ \\
\midrule
\multirow{2}{*}{TPC-C} & default & $634.4$ & $13.3$ & $616$ & $653$ & --- & --- \\
                      & TPB     & $626.2$ & $41.8$ & $553$ & $653$ & $-1.29\%$ & $0.42$ \\
\midrule
\multirow{2}{*}{Rnd-Read}  & default & $15\,552$ & $234$ & $15\,291$ & $15\,877$ & --- & --- \\
                          & TPB     & $15\,815$ & $558$ & $15\,257$ & $16\,727$ & $+1.69\%$ & $0.97$ \\
\midrule
\multirow{2}{*}{Skew-YCSB} & default & $26\,032$ & $476$ & $25\,523$ & $26\,700$ & --- & --- \\
                          & TPB     & $26\,561$ & $373$ & $25\,974$ & $26\,995$ & $+2.03\%$ & $1.95$ \\
\bottomrule
\end{tabular}
\end{table*}

Table~\ref{tab:tpb} summarises the headline result.  TPB matches the
static default within one $\sigma$ on every workload.  On Random Read
and Skewed YCSB the point estimate is \emph{above} the default by
$+1.69\,\%$ and $+2.03\,\%$ respectively; on TPC-C the point estimate
is $-1.29\,\%$, dragged by a single replicate at 553\,k TX/s (the other
four are in the 633--653\,k range, indistinguishable from default).
No comparison reaches significance at $|t|>2.31$.

\textbf{Convergence.}  The final controller state at the end of each
run's last replicate was \texttt{lc=4} (chance $= \tfrac{1}{16}$) on
TPC-C, \texttt{lc=5} (chance $= \tfrac{1}{32}$) on Random Read and on
Skewed YCSB.  These match the sensitivity-sweep optima from \S3
(Table~\ref{tab:sweep}) for TPC-C and Skewed YCSB; Random Read's static
optimum is $\tfrac{1}{16}$ but the controller stayed at the
initialization value of $\tfrac{1}{32}$, suggesting the probe margin
($+0.5\,\%$) was too tight to reliably detect Random Read's modest
optimum.

% --------------------------------------------------------
\subsection{Lessons from the Development Arc (v3$\to$v5)}
% --------------------------------------------------------

Across the development of TPB, we ran four progressively-redesigned
adaptive controllers.  Table~\ref{tab:journey} summarises the headline
delta on TPC-C and Skewed YCSB.

\begin{table}[h]
\centering\small
\caption{Adaptive-controller iterations on the same hardware.  Negative
  $\Delta$ versus the static $\tfrac{1}{32}$ baseline.}
\label{tab:journey}
\begin{tabular}{l rr}
\toprule
Variant & TPC-C $\Delta$ & Skew-YCSB $\Delta$ \\
\midrule
CBP-v3 (per-class hill climber)               & $-11.5\%$ & $-22.4\%$ \\
CBP-v4 (thread-local + hit-rate bandit)       & $-0.83\%$ & $-7.4\%$  \\
TPB-v3 (this $\tau_{\text{phase}}\!=\!15$, $p_e\!=\!8$) & $-2.21\%$ (sig.) & --- \\
\textbf{TPB-v5} (this work, $p_e\!=\!16$)     & $-1.29\%$ & $+2.03\%$ \\
\bottomrule
\end{tabular}
\end{table}

Two architectural lessons emerge:

\textbf{L1.  Don't instrument inside the hot path.}  CBP-v3 added 64
atomic-relaxed \texttt{fetch\_add} operations per access (sampled
1-in-8).  At 25\,M TX/s these alone cost 23\,\% of throughput, before
any adaptation logic runs.  CBP-v4 moved to thread-local counters
flushed in bulk, dropping the overhead to $\approx 7\,\%$ at the same
throughput class.  TPB removes the per-access counters entirely; the
overhead floor is $\approx 0\,\%$.

\textbf{L2.  Pick a signal that aligns with the objective.}
Hit-rate (the fraction of accesses that hit the predicted frame) is
\emph{monotonically increasing} in $p_0$ over almost all workloads.  A
controller that maximises hit-rate therefore drifts toward
$p_0 \to 1$ (always-promote), which is precisely the wrong direction:
the actual throughput optimum sits at small but non-zero $p_0$.  CBP-v4
suffered this failure mode and converged to extreme \texttt{log\_chance}
values across all classes.  TPB uses the system's own throughput
counter, which is exactly the objective function we care about, and
correctly stops at intermediate $p_0$.

% --------------------------------------------------------
\subsection{Robustness of the Result}
% --------------------------------------------------------

Two notes on the strength of the claim.

\textbf{The point estimates lean positive.}  Although none of the
three workloads show statistical significance at $n\!=\!5$, the point
estimates are $+1.69\%$ and $+2.03\%$ on the two read workloads.
This is consistent with TPB correctly identifying the optimal $p_0$
(Skewed YCSB: $\tfrac{1}{32}$ is the optimum, controller stays there)
while paying essentially no overhead.

\textbf{Phase-tolerance and probe-rate sensitivity.}  Table~\ref{tab:knobs}
shows the TPC-C result across a $2\!\times\!2$ grid on
$\tau_{\text{phase}}$ and \texttt{PROBE\_EVERY}.  The dominant knob is
\texttt{PROBE\_EVERY}: halving the probe rate (8$\to$16) reduces the
exploration tax from $\approx 2\,\%$ to within noise.

\begin{table}[h]
\centering\small
\caption{TPC-C grid (n=5, 30\,s).  \texttt{PROBE\_EVERY=16} is
  the dominant knob.}
\label{tab:knobs}
\begin{tabular}{rrr rr}
\toprule
$\tau_{\text{phase}}$ & \texttt{PROBE\_EVERY} & Mean (k) & $\Delta$ & $|t|$ \\
\midrule
5\%   & 8  & 661.8 & $-0.09\%$ & 0.11 \\
100\% & 8  & 660.9 & $-0.22\%$ & 0.23 \\
\textbf{15\%}  & \textbf{16} & \textbf{662.8} & $\bm{+0.07\%}$ & \textbf{0.09} \\
100\% & 16 & 653.1 & $-1.40\%$ & 0.78 \\
\bottomrule
\end{tabular}
\end{table}

% ======================================================================
\section{Related Work}
% ======================================================================

\textbf{Adaptive cache replacement.}
DynamicAdaptiveClimb~\cite{adaptiveclimb2025} adjusts the
\emph{promotion distance} of a CLIMB-style LRU list using a
single tunable based on hit/miss patterns.  This operates on
list-position semantics, not on buffer-pool \emph{frame placement},
which is the substrate TPB targets.  WATT~\cite{watt}
and HyperBolic~\cite{hyperbolic} use sampling-based eviction; both are
orthogonal to promotion-rate adaptation and remain unimplemented for
PrediCache, as flagged in~\cite{zinsmeister2026predictive}.

\textbf{DBMS auto-tuning.}  OtterTune~\cite{ottertune} and follow-ups
use offline ML to recommend hyperparameter values from telemetry.  TPB
solves a much smaller problem (one integer parameter) but does so
\emph{online}, in-process, with negligible compute, which is the
operational regime relevant to a buffer pool.

\textbf{Counter-factual bandits in systems.}  Time-slicing A/B
comparisons are common in scheduler research (e.g.,
RobinHood~\cite{robinhood}).  TPB's contribution is not the
A/B-time-slicing pattern but the observation that it can be hosted
entirely on a system's existing telemetry path, with no new
instrumentation in the data plane.

% ======================================================================
\section{Limitations and Future Work}
% ======================================================================

\textbf{Hardware breadth.}  All measurements are on a single
2-socket Xeon Gold platform.  We have not validated TPB on
AMD EPYC (the platform on which PrediCache's flagship 192-thread
numbers were collected~\cite{zinsmeister2026predictive}), nor on a
direct-attached NVMe.  Lustre's I/O behaviour confines our claims to
the in-memory regime; out-of-memory results would require a real NVMe
device.

\textbf{Run length.}  30-second runs heavily front-load the probe
schedule; only $\approx 1$ probe-sandwich completes before steady
state.  Longer runs (5--30 minutes, typical of production OLTP
benchmarking) would let the controller fully converge and amortise
the exploration tax, likely widening the gap in TPB's favour.

\textbf{Replication.}  $n\!=\!5$ leaves the $\pm 2\%$ point estimates
on Random Read and Skewed YCSB within one $\sigma$ of zero.  A 95\,\%
confidence interval claim would require $\geq 15$ replicates per cell
or longer runs.

\textbf{Signal alternatives.}  TPB's probe-window measurement is
$\textsf{tx\_delta}$, summed across all worker threads.  An
alternative signal --- per-arm \texttt{rdtsc} cycles-per-access ---
would directly measure the cost we care about
(\textsf{cycles}/\textsf{op}) and would generalise to dollar-per-access
in cloud-native costings~\cite{cloud_five_minute_rule}.  We have not
implemented this variant.

\textbf{Generalisation beyond TPB.}  The architectural lesson
(\S5.3 L1, L2) is broader than promotion-probability control.  Any
data-plane parameter with a measurable downstream throughput signal
can be tuned with the same pattern: probe-sandwich on the existing
telemetry tick, zero hot-path overhead.

% ======================================================================
\section{Conclusion}
% ======================================================================

We have presented TPB, an online adaptive controller for the
promotion-probability parameter of PrediCache's predictive-translation
buffer pool.  TPB matches the static $\tfrac{1}{32}$ default within
noise on TPC-C, Random Read, and Skewed YCSB at 16 threads on
commodity hardware, without manual per-workload tuning.  The design
is small (one global mask, three environment knobs, a 50-line
controller hosted by the existing stat thread) and free of
per-access instrumentation, eliminating the 8--23\,\% bookkeeping
floor that doomed our earlier hot-path-instrumented variants.

The most transferable contribution is not TPB itself but the
underlying architectural decision: in a data-plane system, the
controller belongs at the system's coarsest existing telemetry
boundary, not on the hot path.

% ======================================================================
\begin{thebibliography}{99}\small
\bibitem{zinsmeister2026predictive}
M.~Zinsmeister, L.-D.~Nguyen, V.~Leis, T.~Neumann.
\newblock Predictive Translation: High-Performance Buffer Management
  Without the Trade-Offs.
\newblock In \textit{Proc.~SIGMOD}, 2026. DOI: \texttt{10.1145/3786678}.

\bibitem{zinsmeister2026predictive_repo}
PrediCache prototype.
\newblock
  \url{https://github.com/mzinsmeister/PrediCache}. Accessed 2026-05-08.

\bibitem{leis2018leanstore}
V.~Leis, M.~Haubenschild, A.~Kemper, T.~Neumann.
\newblock LeanStore: In-Memory Data Management Beyond Main Memory.
\newblock In \textit{Proc.~ICDE}, 2018.

\bibitem{leis2023vmcache}
V.~Leis, A.~Alhomssi, T.~Ziegler, Y.~Loeck, C.~Dietrich.
\newblock Virtual-Memory Assisted Buffer Management.
\newblock In \textit{Proc.~SIGMOD}, 2023.

\bibitem{adaptiveclimb2025}
Anonymous.
\newblock Adaptive Cache Replacement with Dynamic Resizing.
\newblock arXiv preprint 2511.21235, 2025.

\bibitem{watt}
A.~Alhomssi, V.~Leis.
\newblock Contention and Space Management in B-Trees.
\newblock In \textit{Proc.~CIDR}, 2021.

\bibitem{hyperbolic}
A.~Blankstein, S.~Sen, M.~Freedman.
\newblock Hyperbolic Caching: Flexible Caching for Web Applications.
\newblock In \textit{Proc.~USENIX ATC}, 2017.

\bibitem{ottertune}
D.~Van~Aken, A.~Pavlo, G.~Gordon, B.~Zhang.
\newblock Automatic Database Management System Tuning Through
  Large-Scale Machine Learning.
\newblock In \textit{Proc.~SIGMOD}, 2017.

\bibitem{robinhood}
D.~Berger, B.~Berg, T.~Zhu, S.~Sen, M.~Harchol-Balter.
\newblock RobinHood: Tail Latency Aware Caching --- Dynamic Reallocation
  from Cache-Rich to Cache-Poor.
\newblock In \textit{Proc.~OSDI}, 2018.

\bibitem{cloud_five_minute_rule}
K.~Duwe et al.
\newblock The Five-Minute Rule for the Cloud: Caching in Analytics Systems.
\newblock In \textit{Proc.~CIDR}, 2025.

\end{thebibliography}

\end{document}
